\begin{center}\emph{Allgemeines über Darstellungen von Zahlen durch Formen und ihren Zusammenhang mit Transformationen.}\end{center}

\Article{166}

\emph{Enthält die Form $F$ die Form $F'$, so lässt sich jede durch $F'$ darstellbare Zahl auch durch $F$ darstellen.}

Die Unbestimmten der Formen $F,F'$ seien der Reihe nach $x,y;x',y'$. Nehmen wir an, die Zahl $M$ werde durch $F'$ dargestellt, wenn man $x'=m$, $y'=n$ setzt, und die Form $F$ gehe in $F'$ durch die Substitution

\[
x=\alpha x'+\beta y',\quad y=\gamma x'+\delta y'
\]

\SourceLeaf{147}

über. Dann nimmt offenbar bei der Wahl

\[
x=\alpha m+\beta n,\quad y=\gamma m+\delta n
\]

die Form $F$ den Wert $M$ an.

Lässt sich $M$ auf mehrere Arten durch die Form $F'$ darstellen, zum Beispiel auch durch die Wahl $x'=m'$, $y'=n'$, so folgen daraus mehrere Darstellungen von $M$ durch $F$. Gälte nämlich zugleich

\[
\alpha m+\beta n=\alpha m'+\beta n'\quad\text{und}\quad\gamma m+\delta n=\gamma m'+\delta n'
\]

so wäre entweder $\alpha\delta-\beta\gamma=0$, also gegen die Voraussetzung auch die Determinante der Form $F=0$, oder es wäre $m=m'$, $n=n'$. Daraus folgt, dass $M$ sich durch $F$ auf mindestens ebenso viele verschiedene Arten darstellen lässt wie durch $F'$.

Enthält also sowohl $F$ die Form $F'$ als auch $F'$ die Form $F$, sind also \emph{mit anderen Worten} $F,F'$ äquivalent, und lässt sich die Zahl $M$ durch eine von beiden darstellen, so lässt sie sich auch durch die andere darstellen, und zwar durch die eine auf ebenso viele verschiedene Arten wie durch die andere.

Schließlich bemerken wir, dass in diesem Fall der größte gemeinsame Teiler der Zahlen $m,n$ gleich dem größten gemeinsamen Teiler von $\alpha m+\beta n$, $\gamma m+\delta n$ ist. Der erste sei $=\Delta$, und die Zahlen $\mu,\nu$ seien so gewählt, dass $\mu m+\nu n=\Delta$ gilt. Dann ist

\[
\begin{gathered}(\delta\mu-\gamma\nu)(\alpha m+\beta n)-(\beta\mu-\alpha\nu)(\gamma m+\delta n)\\=(\alpha\delta-\beta\gamma)(\mu m+\nu n)=\pm\Delta\end{gathered}
\]

Also teilt der größte gemeinsame Teiler von $\alpha m+\beta n$, $\gamma m+\delta n$ die Zahl $\Delta$. Umgekehrt teilt $\Delta$ auch diesen Teiler, weil es offenbar $\alpha m+\beta n$, $\gamma m+\delta n$ teilt. Der größte gemeinsame Teiler ist daher notwendig $=\Delta.$ --- Sind also $m,n$ teilerfremd, so sind auch $\alpha m+\beta n.\quad\gamma m+\delta n$ teilerfremd.

\Needspace{12\baselineskip}
\Article{167}

\textsc{Satz.} \emph{Die Formen}

\[
axx+2bxy+cyy\quad .\ .\ .\ .\ .\ .\ (F)
\]

\[
a'x'x'+2b'x'y'+c'y'y'\quad .\ .\ .\ .\ .\ .\ (F')
\]

\emph{seien äquivalent, ihre Determinante $=D$, und die zweite gehe in die erste über, wenn man}

\[
x'=\alpha x+\beta y,\quad y'=\gamma x+\delta y
\]

\SourceLeaf{148}

\emph{setzt. Ferner werde die Zahl $M$ durch $F$ dargestellt, wenn man $x=m$, $y=n$ setzt, und daher durch $F'$, wenn man}

\[
x'=\alpha m+\beta n=m',\quad y'=\gamma m+\delta n=n'
\]

\emph{setzt, und zwar so, dass $m$ zu $n$ und damit auch $m'$ zu $n'$ teilerfremd ist. Dann gehören beide Darstellungen entweder zu demselben Wert oder zu entgegengesetzten Werten des Ausdrucks $\surd D(\operatorname{mod.}M)$, je nachdem, ob die Transformation der Form $F'$ in $F$ eigentlich oder uneigentlich ist.}

\emph{Beweis.} Man bestimme die Zahlen $\mu,\nu$ so, dass $\mu m+\nu n=1$ gilt, und setze

\[
\frac{\delta\mu-\gamma\nu}{\alpha\delta-\beta\gamma}=\mu',\quad \frac{-\beta\mu+\alpha\nu}{\alpha\delta-\beta\gamma}=\nu'
\]

Diese Zahlen sind wegen $\alpha\delta-\gamma\beta=\pm1$ ganzzahlig. Dann gilt

\[
\mu'm'+\nu'n'=1.\quad\text{(Siehe das Ende des vorigen Artikels.)}
\]

Ferner sei

\[
\begin{gathered}\mu(bm+cn)-\nu(am+bn)=V,\\\mu'(b'm'+c'n')-\nu'(a'm'+b'n')=V'\end{gathered}
\]

Dann sind $V,V'$ die Werte des Ausdrucks $\surd D(\operatorname{mod.}M)$, zu denen die erste und die zweite Darstellung gehören. Setzt man in $V'$ für $\mu',\nu',m',n'$ ihre Werte und in $V$

\[
\begin{array}{lll}\text{für}&a,&a'\alpha\alpha+2b'\alpha\gamma+c'\gamma\gamma\\\text{für}&b,&a'a\beta+b'(\alpha\delta+\beta\gamma)+c'\gamma\delta\\\text{für}&c,&a'\beta\beta+2b'\beta\delta+c'\delta\delta\end{array}
\]

ein, so erhält man nach Ausmultiplizieren $V=V'(\alpha\delta-\beta\gamma).$

Also gilt entweder $V=V'$ oder $V=-V'$, je nachdem, ob $\alpha\delta-\beta\gamma=+1$ oder $=-1$ ist. \emph{Das heißt}, die Darstellungen gehören zu demselben Wert oder zu entgegengesetzten Werten des Ausdrucks $\surd D(\operatorname{mod.}M)$, je nachdem, ob die Transformation der Form $F'$ in $F$ eigentlich oder uneigentlich ist. $\text{W.\,z.\,b.\,w.}$

Hat man also mehrere Darstellungen der Zahl $M$ durch die Form $(a,b,c)$ mit teilerfremden Werten der Unbestimmten $x,y$, die zu \emph{verschiedenen} Werten des Ausdrucks $\surd D(\operatorname{mod.}M)$ gehören, so gehören die entsprechenden Darstellungen durch die Form $(a',b',c')$ jeweils zu denselben Werten. Gibt es keine Darstellung der Zahl $M$ durch eine Form, die zu einem bestimmten Wert gehört, so gibt es auch keine Darstellung zu diesem Wert durch eine dazu äquivalente Form.

\SourceLeaf{149}

\Article{168}

\textsc{Satz.} \emph{Wird die Zahl $M$ durch die Form $axx+2bxy+cyy$ dargestellt, indem man $x,y$ die teilerfremden Werte $m,n$ zuweist, und ist der Wert des Ausdrucks $\surd D(\operatorname{mod.}M)$, zu dem diese Darstellung gehört, gleich $N$, so sind die Formen $(a,b,c)$, $(M,N,\frac{NN-D}{M})$ eigentlich äquivalent.}

\emph{Beweis.} Aus Artikel $155$ folgt, dass sich ganze Zahlen $\mu,\nu$ finden lassen, sodass

\[
m\mu+n\nu=1,\quad \mu(bm+cn)-\nu(am+bn)=N
\]

gilt. Damit geht die Form $(a,b,c)$ durch die Substitution $x=mx'-\nu y'$, $y=nx'+\mu y'$, die offenbar eigentlich ist, in eine Form mit der Determinante $=$ $D(m\mu+n\nu)^2$, \emph{das heißt} $=D$, also in eine äquivalente Form über. Schreiben wir diese Form $=(M',N',\frac{N'N'-D}{M'})$, so gilt

\[
\begin{gathered}M'=amm+2bmn+cnn=M,\\N'=-m\nu a+(m\mu-n\nu)b+n\mu c=N\end{gathered}
\]

\begingroup\rightskip=0pt plus 3em\relax
Die Form, in die $(a,b,c)$ durch jene Transformation übergeht, ist also $(M,N,\frac{NN-D}{M})$. $\text{W.\,z.\,b.\,w.}$
\par\endgroup

Aus den Gleichungen

\[
m\mu+n\nu=1,\quad \mu(mb+nc)-\nu(ma+nb)=N
\]

folgt außerdem

\[
\begin{gathered}\mu=\frac{nN+ma+nb}{amm+2bmn+cnn}=\frac{nN+ma+nb}{M},\\\nu=\frac{mb+nc-mN}{M}\end{gathered}
\]

Diese Zahlen sind daher ganzzahlig.

Ferner ist zu beachten, dass die Aussage nicht gilt, wenn $M=0$ ist. Dann wird nämlich der Term $\frac{NN-D}{M}$ unbestimmt\HistoricalNote{*)}{Wollen wir unsere Ausdrucksweise auch auf diesen Fall ausdehnen, so bedeutet die Aussage, $N$ sei ein Wert des Ausdrucks $\surd D(\operatorname{mod.}M)$ oder es gelte $NN\equiv D(\operatorname{mod.}M)$, dass $NN-D$ ein Vielfaches von $M$ und daher $=0$ ist.}.

\Article{169}

Gibt es mehrere Darstellungen der Zahl $M$ durch $(a,b,c)$, die zu demselben Wert des Ausdrucks $\surd D(\operatorname{mod.}M)$, $N$ gehören, wobei wir die Werte von $x,y$ stets als teilerfremd voraussetzen, so lassen sich daraus auch mehrere eigentliche Transformationen der Form $(a,b,c)..(F)$ in $(M.\ N,\frac{NN-D}{M})..(G)$ ableiten. Entsteht nämlich auch aus den Werten $x=m'$, $y=n'$ eine solche Darstellung, so geht $(F)$ auch durch die Substitution

\SourceLeaf{150}

\[
\begin{gathered}x=m'x'+\frac{m'N-m'b-n'c}{M}y',\\y=n'x'+\frac{n'N+m'a+n'b}{M}y'\end{gathered}
\]

in $(G)$ über. Umgekehrt ergibt jede eigentliche Transformation der Form $(F)$ in $(G)$ eine Darstellung der Zahl $M$ durch die Form $(F)$, die zum Wert $N$ gehört. Geht nämlich $(F)$ in $(G)$ über, indem man $x=mx'-\nu y'$, $y=nx'+\mu y'$ setzt, so wird $M$ durch $(F)$ dargestellt, wenn man $x=m$, $y=n$ setzt. Da hier $m\mu+n\nu=1$ gilt, ist der Wert des Ausdrucks $\surd D(\operatorname{mod.}M)$, zu dem die Darstellung gehört, gleich $\mu(bm+cn)-\nu(am+bn)$, \emph{das heißt} $N$. Aus mehreren verschiedenen eigentlichen Transformationen folgen ebenso viele verschiedene Darstellungen, die zu $N$ gehören\HistoricalNote{*)}{Nimmt man an, dass zwei verschiedene eigentliche Transformationen dieselbe Darstellung liefern, so müssen sie folgende Gestalt haben:
\[
\begin{aligned}
1)\quad &x=mx'-\nu y',\quad y=nx'+\mu y';\\
2)\quad &x=mx'-\nu'y',\quad y=nx'+\mu'x'.
\end{aligned}
\]
Aus den beiden Gleichungen
\[
\begin{gathered}
m\mu+n\nu=m\mu'+n\nu',\\
\mu(mb+nc)-\nu(ma+nb)=\mu'(mb+nc)-\nu'(ma+nb),
\end{gathered}
\]
folgt aber leicht, dass entweder $M=0$ oder $\mu=\mu'$, $\nu=\nu'$ gilt. Den Fall $M=0$ haben wir bereits ausgeschlossen.}. --- Kennt man alle eigentlichen Transformationen der Form $(F)$ in $(G)$, so ergeben sich daraus also leicht alle Darstellungen von $M$ durch $(F)$, die zum Wert $N$ gehören. Damit ist die Aufgabe, die Darstellungen einer gegebenen Zahl durch eine gegebene Form zu finden, bei denen die Unbestimmten teilerfremde Werte erhalten, auf die Aufgabe zurückgeführt, alle eigentlichen Transformationen jener Form in eine gegebene äquivalente Form zu finden.

Wenden wir nun das in Artikel $162$ Gezeigte darauf an, so folgt leicht: Ist eine Darstellung der Zahl $M$ durch die Form $(F)$, die zum Wert $N$ gehört, durch $x=\alpha$, $y=\gamma$ gegeben, so lautet die allgemeine Formel, die alle Darstellungen derselben Zahl durch die Form $(F)$, die zum Wert $N$ gehören, umfasst:

\[
x=\frac{\alpha t-(\alpha b+\gamma c)u}{m},\quad y=\frac{\gamma t+(\alpha a+\gamma b)u}{m}
\]

Dabei ist $m$ der größte gemeinsame Teiler der Zahlen $a,2b,c$, und $t,u$ durchlaufen alle Zahlen, die der Gleichung $tt-Duu=mm$ genügen.

\Needspace{14\baselineskip}
\Article{170}

\begingroup\rightskip=0pt plus 3em\relax
Ist die Form $(a,b,c)$ einer ambigen Form äquivalent und damit der Form $(M,N,\frac{NN-D}{M})$ sowohl eigentlich als auch uneigentlich oder sowohl der Form $(M,N,\frac{NN-D}{M})$ als auch der Form $(M,-N,\frac{NN-D}{M})$ eigentlich äquivalent, so gibt es Darstellungen der Zahl $M$ durch die Form \SourceLeaf{151}$(F)$, die sowohl zum Wert $N$ als auch zum Wert $-N$ gehören. Gibt es umgekehrt mehrere Darstellungen der Zahl $M$ durch dieselbe Form $(F)$, die zu \emph{entgegengesetzten} Werten des Ausdrucks $\surd D(\operatorname{mod.}M)$, $N,-N$ gehören, so ist die Form $(F)$ der Form $(G)$ sowohl eigentlich als auch uneigentlich äquivalent, und es lässt sich eine ambige Form angeben, der $(F)$ äquivalent ist.
\par\endgroup

Dies möge hier an allgemeinen Ergebnissen über Darstellungen genügen. Über Darstellungen, bei denen die Unbestimmten keine teilerfremden Werte haben, werden wir weiter unten sprechen. Hinsichtlich der anderen Eigenschaften müssen Formen mit negativer Determinante ganz anders behandelt werden als Formen mit positiver Determinante. Wir werden deshalb nun beide Arten getrennt betrachten und mit den ersten als den leichteren beginnen.
